	\maketitle
	\tableofcontents
	\newpage
	
	% ============================================================
	\section{Introduction: The Operating Point Between Bit and Qubit}
	\label{sec:introduction}
	% ============================================================
	
	Doc.~178 established the stability of spin as macroscopic evidence
	for discrete torsion configurations. The permanent magnet shows:
	spin-up and spin-down are geometrically distinguished states, robust over
	years without cooling. The classical bit uses precisely this regime ---
	a deep potential well $U \gg k_B T$.
	
	The qubit operates at the opposite extreme: metastable torsion,
	coherence in the millisecond range (Doc.~174), enforced by
	extreme cooling to mK temperatures.
	
	The \textbf{p-bit} (probabilistic bit) occupies the conceptually
	hitherto unexplored intermediate range: a magnetic nanostructure
	on the geometric scale $L_8 \approx 22\,\text{nm}$ of the
	$\xi$-hierarchy, at which the torsion barrier becomes thermally
	surmountable.
	
	\begin{significance}[The central thesis of this document]
		The thermal noise of the p-bit is not fundamental randomness.
		It is the macroscopically unresolved sampling of deterministic
		torsion configurations. The apparent stochasticity is epistemic
		--- the transitions follow the sub-Planck-scale geometry
		operating on the timescale $t_0 = t_P/7500$.
		
		The physical barrier is given materially by the
		Stoner-Wohlfarth energy $U = K_u \cdot V$.
		The FFGFT identifies: (a)~\textbf{where} natural p-bit candidates
		are geometrically to be found (scale $L_N$), and (b)~\textbf{how} the
		apparent randomness is to be understood physically.
	\end{significance}
	
	% ============================================================
	\section{The \texorpdfstring{$\xi$}{xi}-Length Hierarchy
		(Review of Doc.~173)}
	\label{sec:length-hierarchy}
	% ============================================================
	
	Doc.~173 defines the geometric scaling ladder of the FFGFT
	via the parameter $\xi_\text{par} = 4/30000 \approx 1{,}333\times10^{-4}$:
	\begin{equation}
		L_N \;=\; L_0 \cdot \left(\frac{1}{\xi_\text{par}}\right)^{\!N},
		\qquad L_0 \;=\; \xi_\text{par}\cdot \ell_P,
		\label{eq:length-hierarchy}
	\end{equation}
	where $\ell_P = 1{,}616\times10^{-35}\,\text{m}$ is the Planck length.
	Level $N=8$ gives:
	\begin{equation}
		L_8 \;=\; \ell_P\cdot\xi_\text{par}^{-7}
		\;\approx\; 21{,}6\,\text{nm}.
		\label{eq:L8}
	\end{equation}
	This scale is identified in Doc.~173 as the universal geometric anchor:
	excitonic Bohr radius, fine structure constant,
	transistor threshold size --- all at $N \approx 8{,}00$.
	
	% ============================================================
	\section{From Length Scale to Characteristic Energy}
	\label{sec:energy-derivation}
	% ============================================================
	
	\subsection{Torsion field characteristic at scale
		\texorpdfstring{$L_N$}{L\_N}}
	\label{subsec:torsion-energy}
	
	Each geometric scale $L_N$ is assigned a characteristic
	torsion energy in the FFGFT framework. The torsion field is a relativistic
	field object; the natural energy scale of a torsion mode at
	wavelength $L_N$ is:
	\begin{equation}
		E_N \;\equiv\; \frac{\hbar c}{L_N}
		\;=\; \frac{\hbar c}{L_0}\,\xi_\text{par}^{N}
		\;=\; \frac{E_P}{\xi_\text{par}}\,\xi_\text{par}^{N}
		\;=\; E_P\cdot\xi_\text{par}^{N-1},
		\label{eq:energy-hierarchy}
	\end{equation}
	where $E_P = \hbar c/\ell_P = \hbar/t_P \approx 1{,}956\times10^9\,\text{J}$
	is the Planck energy. Since $\xi_\text{par} < 1$, $E_N$ decreases strictly
	monotonically with increasing $N$; each level is a factor of
	$1/\xi_\text{par} = 7500$ smaller than the previous one.
	
	\begin{keyresult}[Verification: $N=1$ yields the Planck energy]
		For $N = 1$:
		\begin{equation}
			E_1 \;=\; E_P \cdot \xi_\text{par}^{0} \;=\; E_P
			\;\approx\; 1{,}956\times10^9\,\text{J}.
		\end{equation}
		The Planck energy is the torsion energy of the Planck level --- not a
		free parameter, but a geometric consequence of the hierarchy.
	\end{keyresult}
	
	\subsection{Complete energy hierarchy and the p-bit regime}
	\label{subsec:energy-table}
	
	\begin{center}
		\resizebox{\textwidth}{!}{%
			\begin{tabular}{rllll}
				\toprule
				$N$ & $L_N$ & $E_N = \hbar c / L_N$ &
				$E_N / k_BT_{300}$ & Physical domain \\
				\midrule
				0 & $2{,}16\times10^{-39}$\,m & $1{,}47\times10^{13}$\,J &
				$3{,}5\times10^{33}$ & sub-Planck ($L_0$) \\
				1 & $1{,}62\times10^{-35}$\,m & $1{,}96\times10^{9}$\,J &
				$4{,}7\times10^{29}$ & Planck scale; $E_1 = E_P$ \\
				2 & $1{,}21\times10^{-31}$\,m & $2{,}61\times10^{5}$\,J &
				$6{,}3\times10^{25}$ & --- \\
				3 & $9{,}09\times10^{-28}$\,m & $3{,}48\times10^{1}$\,J &
				$8{,}4\times10^{21}$ & --- \\
				4--7 & \ldots & \ldots & $10^{18}$--$10^{6}$ & --- \\
				\textbf{8} &
				$\mathbf{2{,}16\times10^{-8}}$\,\textbf{m} &
				$\mathbf{1{,}47\times10^{-18}}$\,\textbf{J} &
				$\mathbf{354}$ &
				\textbf{Exciton, Bohr radius, transistor threshold} \\
				\textbf{9} &
				$1{,}62\times10^{-4}$\,m &
				$1{,}95\times10^{-22}$\,J &
				$\mathbf{0{,}047}$ &
				$0{,}16$\,mm \\
				10 & $1{,}21\times10^{0}$\,m & $2{,}61\times10^{-26}$\,J &
				$6\times10^{-6}$ & $\sim 1$\,m \\
				\bottomrule
			\end{tabular}
		}
	\end{center}
	
	\begin{important}[The gap across the p-bit regime]
		The p-bit operating regime $E_N \approx k_B T$ at $T = 300\,\text{K}$
		lies at $N_\text{opt} \approx 8{,}66$ --- \emph{between} the
		integer levels $N=8$ and $N=9$. The step factor is
		$1/\xi_\text{par} = 7500$: between $E_8$ ($354\,k_BT$) and
		$E_9$ ($0{,}047\,k_BT$) there is no integer level value
		that hits $k_BT(300\,\text{K})$.
		
		\textbf{Consequence:} The $\xi$-energy hierarchy $E_N$ does
		\emph{not} directly fix the barrier height of the p-bit. It determines
		the \textbf{geometric scale} of the nanostructure at which
		p-bit-capable resonators naturally occur.
		
		For cryogenic p-bits, however: $E_9 \approx k_BT$ at
		$T^* \approx 14\,\text{K}$. At this temperature a macroscopic
		resonator of dimension $L_9 \approx 0{,}16\,\text{mm}$
		lies exactly on a hierarchy level.
	\end{important}
	
	% ============================================================
	\section{The Physical Barrier: Stoner-Wohlfarth Energy}
	\label{sec:stoner-wohlfarth}
	% ============================================================
	
	\subsection{Derivation}
	\label{subsec:sw-derivation}
	
	A magnetic nanoresonator with uniaxial anisotropy possesses an
	energy barrier between spin-up and spin-down, given in the
	Stoner-Wohlfarth model \cite{Stoner1948} by:
	\begin{equation}
		U_\text{sw} \;=\; K_u \cdot V
		\label{eq:stoner-wohlfarth}
	\end{equation}
	Here $K_u$ is the uniaxial anisotropy constant
	(in $\text{J/m}^3$, material-specific) and $V$ is the volume of the
	magnetic free layer.
	
	In the FFGFT framework, $U_\text{sw}$ is the macroscopic manifestation
	of the torsion barrier between the two stable configurations
	that Doc.~178 identifies as geometrically distinguished states.
	
	For a spherical nanoresonator with diameter $D$:
	\begin{equation}
		U_\text{sw}(D) \;=\; K_u\,\cdot\,\frac{\pi}{6}\,D^3.
		\label{eq:barrier-geometric}
	\end{equation}
	The condition for p-bit operation $U_\text{sw} \approx k_B T$ fixes
	a specific diameter $D^*$ for a given $K_u$:
	\begin{equation}
		D^* \;=\; \left(\frac{6\,k_B T}{\pi\,K_u}\right)^{1/3}.
		\label{eq:d-star}
	\end{equation}
	
	\subsection{Numerical values and comparison with \texorpdfstring{$L_8$}{L8}}
	\label{subsec:pbit-regime}
	
	For superparamagnetic nanoparticles with
	$K_u \approx 10^4\,\text{J/m}^3$:
	
	\begin{center}
		\begin{tabular}{rrrr}
			\toprule
			$D$ (nm) & $U_\text{sw}$ (meV) & $U_\text{sw}/k_BT_{300}$ & Regime \\
			\midrule
			5  & 4{,}1  & 0{,}16 & \\
			8  & 16{,}7 & 0{,}65 & p-bit \\
			10 & 32{,}7 & 1{,}26 & p-bit \\
			12 & 56{,}5 & 2{,}18 & p-bit \\
			15 & 110    & 4{,}3  & p-bit \\
			20 & 261    & 10{,}1 & p-bit \\
			25 & 511    & 19{,}8 & \\
			\bottomrule
		\end{tabular}
	\end{center}
	
	The p-bit operating range ($U/k_BT \approx 1$--$10$) lies at
	$D \approx 8$--$20\,\text{nm}$ --- exactly in the range of the
	$\xi$-hierarchy level $N = 8$ with $L_8 \approx 22\,\text{nm}$.
	
	\begin{foundation}[Geometric prediction: $N=8$ as natural p-bit anchor]
		The FFGFT predicts: magnetic nanoresonators on the
		$\xi$-level $N \approx 8$ are natural p-bit candidates at
		room temperature. Not because $E_8 \approx k_BT$ (that is not the case:
		$E_8 = 354\,k_BT$), but because $L_8 \approx 22\,\text{nm}$ is the
		geometric scale at which the Stoner-Wohlfarth volume with
		typical material values falls into the p-bit regime.
		
		This coincidence is not accidental in the FFGFT:
		$N=8$ is the universal scale of magnetic single resonators
		(Doc.~173: minimal resonator, Doc.~178: permanent magnet as
		collective resonator at $N\approx8$).
	\end{foundation}
	
	% ============================================================
	\section{The Kramers Switching Rate in the FFGFT Framework}
	\label{sec:kramers}
	% ============================================================
	
	\subsection{Standard formula}
	\label{subsec:kramers-standard}
	
	The thermally activated switching rate between torsion configurations
	follows the Kramers formula \cite{Kramers1940, Hanggi1990}:
	\begin{equation}
		\Gamma \;=\; \frac{\omega_0\,\omega_b}{2\pi\gamma}
		\,\exp\!\left(-\frac{U_\text{sw}}{k_B T}\right),
		\label{eq:kramers}
	\end{equation}
	where $\omega_0$ is the angular frequency at the potential minimum,
	$\omega_b$ at the barrier peak and $\gamma$ is the (Gilbert) damping rate.
	
	\subsection{FFGFT interpretation: determinism instead of randomness}
	\label{subsec:kramers-ffgft}
	
	In the standard interpretation, $\Gamma$ is a stochastic rate:
	each transition is a random, uncorrelated event.
	
	\begin{foundation}[FFGFT reinterpretation of the Kramers rate]
		In the FFGFT, equation~\eqref{eq:kramers} is \textbf{formally identical}
		--- but its physical meaning changes:
		
		$\Gamma$ is the \textbf{time-averaged frequency of deterministic
			torsion transitions}. Each individual transition follows a precise
		path in the sub-Planck-scale geometry of the torsion field,
		which is determined on the timescale $t_0 = t_P/7500$.
		
		The macroscopic measurement --- temporal resolution $\Delta t \gg t_0$ ---
		averages over $\Delta t/t_0 \gg 1$ microstates and creates
		the appearance of randomness. The randomness is \textbf{epistemic},
		not ontological.
		
		This is the direct transfer of the FFGFT position established in
		Doc.~178 (epistemic ignorance rather than genuine superposition)
		to the thermally fluctuating p-bit system.
	\end{foundation}
	
	\subsection{Why no modified prefactor}
	\label{subsec:no-prefactor}
	
	A numerical FFGFT correction to the Kramers rate would require the
	structural parameter $k_BT\,t_0/\hbar$:
	\begin{equation}
		\frac{k_B T\,t_0}{\hbar}
		\;=\; \frac{k_B \cdot 300\,\text{K}}{\hbar/t_0}
		\;\approx\; 2{,}82 \times 10^{-34}.
		\label{eq:structural-parameter}
	\end{equation}
	This value is $\ll 1$ at all technically accessible temperatures.
	A correction to the Kramers prefactor by the FFGFT is not
	experimentally detectable at room temperature.
	
	\begin{keypoint}[Where the FFGFT contribution really lies]
		\begin{enumerate}
			\item \textbf{Geometric scale:} $D \sim L_8 \approx 22\,\text{nm}$
			as natural p-bit anchor --- verifiable materials-science
			prediction.
			\item \textbf{Interpretation:} The switching statistics are deterministically
			structured on the $t_0$ scale; the macroscopically observed
			stochasticity is a projection.
			\item \textbf{No new terms:} The Kramers rate~\eqref{eq:kramers}
			remains unchanged. The FFGFT contribution is interpretive.
		\end{enumerate}
	\end{keypoint}
	
	% ============================================================
	\section{The Three-Regime Picture: Bit, p-Bit, Qubit}
	\label{sec:three-regimes}
	% ============================================================
	
	{\footnotesize
	\begin{center}
		\begin{tabular}{llll}
			\toprule
			\textbf{Type} & \textbf{Barrier} & \textbf{Temp.} & \textbf{Timescale} \\
			\midrule
			Class.\ bit & $U \gg k_BT$ & 300\,K & Years \\[2pt]
			\textbf{p-bit} & $U \approx k_BT$ & 300\,K & ns--$\mu$s \\[2pt]
			Qubit & $U \gg k_BT$ (prep.) & mK & $\mu$s--ms \\
			\bottomrule
		\end{tabular}
	\end{center}
	\begin{center}
		\begin{tabular}{lll}
			\toprule
			\textbf{Type} & \textbf{$\xi$-level} & \textbf{FFGFT: random?} \\
			\midrule
			Class.\ bit & $N \approx 8$, deep & No; stable torsion well \\[2pt]
			\textbf{p-bit} & $D \sim L_8 \approx 22$\,nm & No; deterministic \\[2pt]
			Qubit & $N \approx 8$, metastable & No; epistemic (Doc.~178) \\
			\bottomrule
		\end{tabular}
	\end{center}
	}
	
	All three types reside on the same geometric scale $N \approx 8$.
	The difference is the operating point on the potential curve:
	deep in the minimum (bit), at the tipping point (p-bit), metastably prepared (qubit).
	
	% ============================================================
	\section{Landauer's Principle}
	\label{sec:landauer}
	% ============================================================
	
	Landauer's principle \cite{Landauer1961} applies to the p-bit unchanged:
	\begin{equation}
		E_\text{diss}^\text{min} \;=\; k_B T \ln 2.
		\label{eq:landauer}
	\end{equation}
	The p-bit switches between two torsion configurations; the
	state space is binary. The $\xi$-geometry determines the \emph{dynamics}
	(scale, switching rate), not the \emph{state space size}.
	
	\begin{important}[No modified Landauer bound in the FFGFT]
		A $\xi$-dependent Landauer bound would require the
		state space of the p-bit to have more than two elements, or that
		information erasure uses geometrically non-equivalent paths.
		Neither is the case in a standard p-bit.
		The standard bound $k_BT \ln 2$ applies.
	\end{important}
	
	% ============================================================
	\section{Experimental Realisations}
	\label{sec:experiment}
	% ============================================================
	
	\begin{itemize}
		\item \textbf{Magnetic Tunnel Junctions (MTJ) \cite{Borders2019}:}
		CoFeB free layers, $K_u \sim 10^4$--$10^5\,\text{J/m}^3$,
		lateral dimension $\sim 20$--$100\,\text{nm}$ --- directly at $L_8$.
		Switching times in the nanosecond range.
		
		\item \textbf{Superparamagnetic nanoparticles:}
		Fe$_3$O$_4$, $\gamma$-Fe$_2$O$_3$ with $D \approx 10$--$20\,\text{nm}$
		show spontaneous thermal switching at 300\,K in the
		$L_8$ regime.
		
		\item \textbf{CMOS-based p-bits \cite{Aadit2022}:}
		Ring oscillators near metastability. No direct
		torsion field connection --- realise the logical operating point
		by other means.
	\end{itemize}
	
	\begin{keypoint}[Verifiable FFGFT prediction]
		The FFGFT predicts a \textbf{geometrically distinguished dimension}
		$D^* \approx L_8 \approx 22\,\text{nm}$ for optimal p-bit quality
		--- material-independent, as long as $K_u$ scales accordingly.
		A systematic measurement of switching rate, noise colour and
		reproducibility across different nanoparticle sizes should
		show an optimum at $D \approx L_8$.
	\end{keypoint}
	
	% ============================================================
	\section{Classification within the FFGFT Document Series}
	\label{sec:classification}
	% ============================================================
	
	\begin{center}
	{\footnotesize
	\begin{tabularx}{\textwidth}{clX}
		\toprule
		\textbf{Doc.} & \textbf{Topic} & \textbf{Relation to Doc.~184} \\
			\midrule
			173 & $\xi$-hierarchy, quantum dot, minimal resonator &
			Foundation: $L_N$, $E_N$ derivation, $N=8$ anchor \\
			174 & Reading and writing spin &
			Techniques; decoherence times \\
			178 & Spin stability, permanent magnet &
			Epistemic randomness; torsion well vs.\ tipping point \\
		\textbf{184} & \textbf{p-bit} &
		\textbf{Tipping point: same scale $N=8$, different operating point} \\
		\bottomrule
	\end{tabularx}
	} % footnotesize
\end{center}
	
	% ============================================================
	\section{Summary}
	\label{sec:summary}
	% ============================================================
	
	\begin{conclusionBox}[Conclusion: p-bit in the FFGFT framework]
		\textbf{1. The $\xi$-energy hierarchy is cleanly derivable.}\\
		$E_N = \hbar c/L_N = E_P\cdot\xi_\text{par}^{N-1}$, with $E_1 = E_P$.
		The levels jump by a factor of $7500$ --- the p-bit regime at
		300\,K lies between $N=8$ and $N=9$.
		
		\textbf{2. The barrier is Stoner-Wohlfarth: $U = K_u \cdot V$.}\\
		Adjustable by materials engineering, not quantised by $E_N$.
		p-bit regime at $D \approx 8$--$20\,\text{nm}$ ($K_u \sim 10^4\,\text{J/m}^3$).
		
		\textbf{3. The $\xi$-reference is the geometric scale $L_8 \approx 22\,\text{nm}$.}\\
		This is the universal resonator scale $N=8$ from Doc.~173.
		Natural p-bit candidates lie on this level.
		
		\textbf{4. The Kramers rate remains formally unchanged.}\\
		FFGFT contribution: determinism of transitions on the $t_0$ scale ---
		the stochasticity is epistemic, not ontological.
		
		\textbf{5. The Landauer bound applies unchanged.}\\
		$E_\text{diss}^\text{min} = k_BT\ln 2$ --- binary state space,
		no $\xi$ modification.
		
		\vspace{0.5em}
		\textbf{The p-bit is the geometric tipping point between permanent magnet
			(Doc.~178) and qubit (Doc.~174) --- on the same $\xi$-level $N=8$.}
	\end{conclusionBox}
	
	\newpage
	\begin{thebibliography}{99}
		
		\bibitem{Datta2018}
		Datta, S.\ et al.
		\textit{p-Bits for Probabilistic Spin Logic.}
		arXiv:1809.04028 (2018).
		
		\bibitem{Borders2019}
		Borders, W.\,A.\ et al.
		\textit{Integer factorization using stochastic magnetic tunnel junctions.}
		Nature \textbf{573}, 390--393 (2019).
		
		\bibitem{Aadit2022}
		Aadit, N.\,A.\ et al.
		\textit{Massively parallel probabilistic computing with sparse inverse
			Ising problems.}
		Nature Electronics \textbf{5}, 460--468 (2022).
		
		\bibitem{Stoner1948}
		Stoner, E.\,C.\ \& Wohlfarth, E.\,P.
		\textit{A mechanism of magnetic hysteresis in heterogeneous alloys.}
		Phil.\ Trans.\ R.\ Soc.\ London A \textbf{240}, 599--642 (1948).
		
		\bibitem{Kramers1940}
		Kramers, H.\,A.
		\textit{Brownian motion in a field of force and the diffusion model
			of chemical reactions.}
		Physica \textbf{7}, 284--304 (1940).
		
		\bibitem{Hanggi1990}
		Hänggi, P., Talkner, P.\ \& Borkovec, M.
		\textit{Reaction-rate theory: fifty years after Kramers.}
		Rev.\ Mod.\ Phys.\ \textbf{62}, 251--341 (1990).
		
		\bibitem{Landauer1961}
		Landauer, R.
		\textit{Irreversibility and heat generation in the computing process.}
		IBM J.\ Res.\ Dev.\ \textbf{5}, 183--191 (1961).
		
		\bibitem{Hayakawa2021}
		Hayakawa, K.\ et al.
		\textit{Nanosecond random telegraph noise in in-plane magnetic
			tunnel junctions.}
		Phys.\ Rev.\ Lett.\ \textbf{126}, 117202 (2021).
		
		\bibitem{Pascher2026_173}
		Pascher, J.
		\textit{The Minimal Quantum Resonator.}
		Doc.~173, FFGFT framework (2026).
		
		\bibitem{Pascher2026_174}
		Pascher, J.
		\textit{Reading and Writing Spin.}
		Doc.~174, FFGFT framework (2026).
		
		\bibitem{Pascher2026_178}
		Pascher, J.
		\textit{Stability as Ground State.}
		Doc.~178, FFGFT framework (2026).
		
	\end{thebibliography}
